\chapter{Discussion and Conclusion}

\section{The question, revisited}

This dissertation asked one question with two halves: \textit{can subtractive attention guidance be
delivered on consumer XR hardware, and when it is delivered, does it actually help people focus?}

The first half is answered in the affirmative, with precision about what delivered means. On the Meta
Quest 3 the question reduces to a compositing-rights problem, and the access an application has
stratifies into three tiers: colour-only styling of the OS-owned passthrough layer, overlay
compositing on top of it, and full pixel control over a lower-quality camera feed. That narrow space
turns out to be sufficient for a working multi-mode system, built around an architectural inversion
in which the focus region is rendered as an \textit{absence} in an overlay, so native passthrough at
full quality serves the region that matters while the treatments occupy the periphery around it. The
hybrid composite at the heart of the most capable mode survived a structured prior-art search with no
published or shipped precedent found. The architecture's two central trade-offs were then measured on
the device itself: the camera-free dark modes cost nothing measurable over an empty passthrough scene, and the
native path the window preserves resolves at least two logMAR steps finer than the camera re-render
it spares the user from (Chapter 4, Section 4.11).

The second half is answered by the study, completed at twenty people tested: the pre-registered
twenty through the full two-block protocol, giving nineteen
workstation pairs and seventeen driving pairs, with the decision table applied
as committed. The answers:
a 3.07-point accuracy gain on a real-hardware task, real, significant, and at close to the planned
sensitivity, but
smaller than the pre-set half-the-errors bar, so claimed as a bounded benefit under the table's
fourth row rather than at the pre-set size; peripheral awareness with no measured loss, the safety
hypothesis's refuting observation never occurring and the pooled measure instead rising, though the
two-sided equivalence statistic itself was not established, for the one-directional reason set out in
Chapter 5, Section 5.7.2; and a detection benefit appearing in
the 20--30$^{\circ}$ band the system's own detection-gated geometry predicts, supported at 95\%
power and bounded by the band-assignment proxy of Chapter 5, Section 5.7.3. Every outcome was mapped to a pre-committed conclusion before the data existed, and the
outcome that did occur is
a specific, mechanistically coherent one: the system helps, by a real amount whose exact size a
larger sample would pin more tightly.

RQ3's answer does not simply echo the behavioural one. Perceived focus benefit is strongly positive,
but perceived awareness cost splits close to evenly, with roughly four in ten responses reporting
feeling cut off from the surroundings or noticing side events late (Chapter 5, Section 5.7.5). The
behavioural hit rate shows no such loss, so the subjective impression and the measured outcome
disagree, and that disagreement is reported rather than resolved.

\section{What the design space contributes, independently of the data}

The central finding of Part T does not await any result. It is the shape of the design space. On a
consumer VST headset, diminished reality is not one problem but three, indexed by compositing access.
At Tier 1 the platform grants global colour remapping and nothing spatial: no per-pixel masks, no
world-locked regions, no read access, and, since the deprecation of surface-projected passthrough at
SDK v83, no per-surface control either. The significant fact is the \textit{absence}: a developer who
wants to dim the periphery and spare a window cannot do it at Tier 1 at all. That absence is a
platform-policy finding rather than an engineering one, and it arguably shapes what DR research is
possible on consumer hardware more than any algorithm does.

The system's response is an existence proof that windowed subtractive DR is achievable at Tier 2
without OS cooperation, and that hybrid Tier-2/Tier-3 composites are achievable with Passthrough
Camera Access. The composite, precisely scoped, is the novelty claim, and the documented search
protocol is what entitles the dissertation to make it.

Four engineering lessons generalise beyond this project, and Chapter 4, Section 4.10 records each
with the failure that produced it. \textbf{Sample by world direction, not screen space}, when
re-rendering a mono camera feed, because only an eye-independent sampling domain fuses binocularly.
\textbf{Exploit the quality asymmetry}: on any platform where the accessible feed is worse than the
native view, put the manipulation where the eyes are not, so the task-relevant region never pays the
quality tax. \textbf{Decouple effects from head motion}, inverting the common default of effects that
follow the head into one that yields to it, because head-coupled restriction increases discomfort
\citep{norouzi2018}. And \textbf{consumer-platform DR research inherits platform-policy risk}:
surface-projected passthrough, a capability this project might have built on, was deprecated
mid-project. Tier boundaries on consumer hardware are set by vendor policy and move without notice,
so designs that keep a working path at every tier are more robust to that movement than designs
premised on a single capability.

\section{Limitations}

The limitations below are collected from Chapters 3 to 5 so they appear once, together, rather than
diluted across the text.

\begin{enumerate}
\item \textbf{Mode $\times$ scenario confound.} Hard Dark is tested only in the workstation context
  and SignPop only in the monitoring context, because the three-mode sampler was cut from the running
  protocol. Every empirical claim is mode-in-context. The design cannot separate mode effects from
  scenario effects, and nothing in it mitigates this.
\item \textbf{SignPop's manipulation is a bundle.} A confirmed detection switches on six effects at
  once: colour pop, window carve-out, saturation lift, brightness lift, contrast expansion and
  darkened surround (Chapter 4, Section 4.5.5). None is isolable in the current design, so no
  measured benefit can be attributed to any one of them.
\item \textbf{Soft Dark carries no participant data at all.} It was never placed in front of a
  participant. Its place in the dissertation rests on design-space grounds and on its shader path,
  never on evidence about how it is experienced.
\item \textbf{The oracle gap was never measured.} SignPop is detector-dependent, and the detection
  driving it is an offline full-resolution bake rather than a live on-device pipeline. How far a live
  detector would fall short of that oracle, in recall, throughput and latency, is unquantified. Every
  detector-conditioned statement inherits that qualifier.
\item \textbf{The technical measurements cover two properties, not five.} Frame cost and legibility
  are measured (Chapter 4, Section 4.11); end-to-end latency, world-locking stability, and thermal
  endurance remain argued from structure rather than measured. Within what was measured, one
  criterion was missed: the video scene's stale-frame rate exceeds the pre-stated 1\% bound in both
  study arms equally, a property of the shared 4K playback path rather than of any treatment. The
  legibility read is the author's, from a cast recording, not a psychophysical procedure.
\item \textbf{No eye tracking.} Quest 3 has none. Head pose is a validated proxy at the
  eccentricities used \citep{higgins2022,sitzmann2018}, but covert attention shifts, distraction
  without head movement, are invisible to telemetry. The dependent variables were chosen so that
  \textit{performance outcomes} rather than attention allocation carry the hypotheses, so the
  mechanism remains partially inferred.
\item \textbf{Sample, power, and Block B's response-criterion qualification.} The pre-registered
  twenty completed
  the full two-block protocol, and the primary hypothesis and the 20--30$^{\circ}$ band sit at 88\%
  and 89.0\% achieved power, two-tailed as registered against the design's 80\% target. Block B lost three
  pairs to its response-validity and comprehension exclusions, all reported in Chapter 5 and
  including one discretionary case, made by experimenter decision on the day and disclosed as such.
  The pooled Block B
  estimate carries a response-criterion qualification and sits at 54.0\% achieved power (Chapter 5,
  Section 5.7.2), and it is the one headline estimate that does not survive leave-one-out, holding
  significance in seven of seventeen such samples where the primary contrast and the pre-registered
  band survive every single-participant removal (Chapter 5, Section 5.7.4).
\item \textbf{Band membership is assigned by a proxy.} The eccentricity bands carrying H2a are
  assigned from each target's mid-lifetime position rather than its position at the instant of the
  press, because the press-time reconstruction is not implemented in the analysis pipeline. Targets
  move while visible, the median lifetime swing is 17.2$^{\circ}$, and roughly 31\% of catches land
  in a different band from the one the mid position assigns (Chapter 5, Sections 5.5 and 5.7.3).
  The band result is the dissertation's largest single effect, and this is the measurement caveat
  that bounds it.
\item \textbf{Workload is not among the measured constructs.} Subjective response was collected
  with a purpose-built end-of-session questionnaire covering perceived focus benefit, perceived
  awareness cost and visual comfort, the first two written as the subjective counterparts of the two
  blocks' own questions (Chapter 5, Sections 5.6 and 5.7.5). Task demand in the
  NASA-TLX sense is not one of those constructs and was not collected separately per condition, so
  the study reports what participants thought of the effects but makes no claim about whether the
  vignette made the task feel less effortful as well as more accurate. Separately, the distractor
  reel's cost is taken from the published literature \citep{ai2025} rather than verified internally.
\item \textbf{Single experimenter, single site, single hardware generation, first exposure.} Scripts
  and checklists bound experimenter variance but cannot eliminate it, and every number is specific to
  Quest 3 and the frozen OS and SDK versions. Practice-to-criterion and counterbalancing distribute
  first-exposure effects but do not abolish the fact that every participant meets peripheral
  diminishment for the first time that hour. The workstation distractors are video panels on the
  task display and the monitoring stimulus is a video: the study measures attention mechanisms under
  controlled distraction, deliberately, and its conclusions end where those controls end.
\item \textbf{No demographic data.} The post-session questionnaire carries no age, gender, vision or
  background items, so participant demographics cannot be reported. What is reported instead is the
  sampling frame: adults recruited by direct approach from the postgraduate computer-science cohort,
  unpaid, under the inclusion criteria of Chapter 5, Section 5.3. The frame is narrow, and findings
  generalise from it accordingly.
\item \textbf{The detection system reads position and size, never motion.} A detection aperture
  responds to an object's angular position and apparent size per frame, with no velocity or
  trajectory term, so an approaching pedestrian and a stationary one subtending the same angle are
  treated identically. Approach speed is exactly the cue a driving-relevant deployment would need.
  Relatedly, the aperture budget (eight live, sixteen in the video path) fills in arrival order with
  no prioritisation, so a scene dense enough to exhaust it would degrade by dropping the newest
  detection rather than the least important one. Neither cap was reached in any session, a property
  of the footage rather than a guarantee.
\end{enumerate}

\section{Transfer: what leaves this dissertation, and under what conditions}

\textbf{To optical see-through hardware, the dark overlay does not transfer.} The Tier-2 modes work
because a VST display is an opaque screen under full render control, and additive OST optics cannot
darken the world. On OST, ``draw black'' means ``draw nothing''. What transfers is the taxonomy and the
question. On OST hardware the taxonomy collapses toward display-side attenuation, of which segmented
dimming hardware such as the Magic Leap 2's dedicated low-resolution dimming panel \citep{magic2024}
is the nearest existing analogue of a Tier-2 windowed dim, and toward the additive salience
modulation demonstrated by \citet{sutton2022}. Critically, the \textbf{human-factors findings are
expected to transfer}, because what peripheral diminishment does to attention should depend on the
human visual system more than on the display that implements it. The engineering is platform-bound.
The psychology should be far less so, though display properties such as field of view and latency
could still moderate the effect. A replication of the Chapter 5 paradigm on a dimming-capable OST
device would be the cleanest test of that expectation.

\textbf{From the simulation, transfer is conditional, and the conditions are enumerable.} Monoscopic
video removes motion-in-depth, so peripheral capture by looming stimuli is untested. The passenger
viewpoint removes task coupling, which is why no statement here concerns driving performance. Video
dynamic range compresses glare, so SignPop's glare-suppression track is exercised only weakly. And
identical stimuli for all participants trade ecological breadth for internal validity, a trade that
is intended. The workstation block, running on the real system against physical distractors, is the
one component that transfers without these caveats, which is why it carries the primary hypothesis.

The two blocks also carry an incidental methodological payload. They run the same guidance principle
at opposite ends of the feasibility spectrum, so if they agree in direction the simulation
methodology that the nearest published work relies on entirely (\citealp{cheng2022};
\citealp{mclaughlin2025}, both VR-simulated) gains external credibility from this project's
real-hardware anchor, and if they diverge the divergence localises what simulation misses. Agreement
is read at the level of sign rather than magnitude, since the blocks use different tasks, dependent
variables and modes, so any cross-block statement is observational rather than powered.

\section{Future work}

Six directions follow directly from the material, in the order their absence limits the present
claims.

\textbf{Separating the bundle.} The 20--30$^{\circ}$ finding is the clearest reason this item leads.
A three-arm design, no filter against filter with the detection gate off against filter with it on,
would need no new code, since the gate already has a toggle. It would directly separate ``dimming
helped'' from ``passing the signal through helped''.

\textbf{Closing the oracle gap, and finishing the measurement campaign.} SignPop is
detector-dependent on offline oracle-quality detection. The natural next iteration replaces the
offline gate with a live detector and re-runs the driving block to test whether the measured benefit
survives the accuracy and coverage gap. That programme also needs the measurements this dissertation
did not take: end-to-end latency, world-locking stability, and above all thermal endurance, since the
frame-cost results of Chapter 4, Section 4.11 show the Tier-3 path holding frame rate comfortably
while saying nothing about holding it for an hour. A prioritised aperture budget belongs to the same
iteration: ranking detection candidates by urgency, for example eccentricity, apparent size and
time-to-contact, and spending the fixed slots on the top few, so a crowded scene degrades by
dropping the least critical detection rather than the newest. There may be no clean solution at high
density, since a filter that opens twenty holes has stopped filtering, and at some density the
honest behaviour is to disable the effect rather than dilute it.

\textbf{Dose--response.} The study tests only the extreme of the attenuation axis. Soft Dark, the
midpoint, was built but never tested, so the middle of the axis has no data behind it. A
discrete-levels design, off against partial against full attenuation, counterbalanced, would
establish the shape of the benefit curve and answer the deployment question this dissertation cannot:
how much diminishment buys how much focus, at what comfort cost.

Three further directions follow less directly but are seeded by the existing architecture.
\textbf{Adaptive introduction}: rather than switching the effect on, a deployed system might
introduce it gradually as focus wanes, a question the formation animation and motion-coupled fade
mechanics already make tractable, though as a validation instrument a ramp confounds time with
intensity and was excluded here. \textbf{Hour-long deployment}: the camera-free dark modes are the
only plausible basis for session-length use today, and a longitudinal study with workload and comfort
instruments administered across days is the bridge between controlled findings and any claim about
practice. \textbf{Saliency-weighted diminishment}: the review identified dimming weighted by computed
scene saliency as unstudied in AR attention guidance, and since the shader already computes
per-fragment treatment, driving its intensity from a saliency model of the periphery is a natural
composition of the Tier-3 path with the saliency-modulation literature.

\section{Closing}

The premise of this work was that a consumer headset can be an instrument for \textit{removing} the
world as well as adding to it, and that both halves of that proposition deserve rigour: the systems
half, where an OS-owned display layer had to be subtracted from by purely additive means, and the
human half, where ``it helps you focus'' had to be converted from a hope into a hypothesis that could
fail.

The evaluation apparatus built to do that is part of the contribution. A single primary endpoint with
a smallest effect size of interest, equivalence margins for safety claims, a gated pilot whose
failure aborts rather than degrades the main study, and a decision table committed before data
collection, is not yet standard practice at masters-project scale. The apparatus costs little, since its
artefacts are documents and thresholds rather than equipment, and it converts every outcome into a
claim with stated bounds. One of those bounds turned out to bind on the study itself: the safety
hypothesis was written as a two-sided equivalence test when the question it encoded had only one
side, and Chapter 5 reports both the pre-registered statistic and the mismatch rather than quietly
substituting the better-fitting one.

The system exists and runs on target hardware. The hypothesis was tested, on nineteen workstation
pairs and seventeen driving pairs, and answered under the rules committed to before the data existed:
a real focus benefit, claimed at its bounded size rather than its hoped-for one; peripheral
awareness with no measured loss; the benefit appearing exactly where the geometry said to look. The question was
asked on the hardware people actually own, and that is the standard the answer inherits.
